% (C) 2026 Ralf Stephan, in collaboration with Claude Code.
% Released under CC0 1.0 Universal (public-domain dedication).
% See https://creativecommons.org/publicdomain/zero/1.0/

\documentclass[11pt]{article}
\usepackage[T1]{fontenc}
\usepackage[utf8]{inputenc}
\usepackage{lmodern}
\usepackage[margin=2.5cm]{geometry}
\usepackage{amsmath,amssymb,amsthm}
\usepackage[hidelinks]{hyperref}

% Lean names and paths: verbatim, breakable at . _ /
\DeclareUrlCommand\lean{\urlstyle{tt}}
\def\UrlBreaks{\do\.\do\_\do\/}
\emergencystretch=6em

\newtheorem{theorem}{Theorem}[section]
\newtheorem{corollary}[theorem]{Corollary}
\newtheorem{lemma}[theorem]{Lemma}
\theoremstyle{definition}
\newtheorem{definition}[theorem]{Definition}
\newtheorem{remark}[theorem]{Remark}

\title{Further results on the brief\\ \emph{Physical vector spaces for quantum mechanics}}
\author{Ralf Stephan}
\date{14 November 2026}

\begin{document}
\maketitle

\begin{abstract}
  This note collects results that go beyond the technical brief \emph{Physical vector spaces for
  quantum mechanics}~\cite{brief} and partial answers to its open problems. They are not part of
  the brief. Each section is self-contained, and each names the Lean declarations that check it;
  the Lean files are in the brief's directory, under \lean{QuantumVectorSpaces/FurtherResults/}.
\end{abstract}

\section{The Hilbert brief's change of frame is not a map of Schwartz space}
\label{sec:schwartz}

The brief proposes quantum vector spaces (QSTO:V) as a replacement for the Hilbert-space
formulation QSTO:H, and argues in the proof of its Theorem~5 that QSTO:V satisfies [B]WDQ. The
argument assumes that ``the defining observables of A must be observables in B and
vice-versa'' for two frames $A$ and $B$ of the same system. This premise holds exactly when the
two frames give the same topology on $V$ (checked with the brief,
\lean{QuantumVectorSpaces.topology_eq_iff_continuous}). It is not automatic. This section shows
that it fails for the brief's own Example~3, $V = \mathcal S(\mathbb R)$, with the change of
frame that the companion brief \emph{Unphysicality of Hilbert spaces for quantum
mechanics}~\cite{hilbert} uses against QSTO:H.

That change of frame is $y = \tan(\frac\pi2\operatorname{erf}x)$. Its unitary $U$ on
$L^2(\mathbb R)$ sends the Gaussian packet
\[
  \psi(x) = \bigl(e^{-x^2}/\sqrt\pi\bigr)^{1/2} = \pi^{-1/4}e^{-x^2/2}
  \quad\text{to}\quad
  \phi(y) = \bigl(\pi(1+y^2)\bigr)^{-1/2}.
\]

\begin{theorem}\label{thm:schwartz}
  \begin{enumerate}
    \item $\psi$ is a Schwartz function, and $U\psi = \phi$ is not. So $U$ does not map
      $\mathcal S(\mathbb R)$ into itself.
    \item The position $Y$ of the second frame, carried back to the first as $U^{-1}YU$, is not
      defined at the Schwartz function $\psi$.
  \end{enumerate}
\end{theorem}

\begin{proof}
  Every Schwartz function $f$ lies in the domain of the position operator, since $xf$ is again
  a Schwartz function and so in $L^2$. But $\int y^2|\phi(y)|^2\,dy =
  \int y^2/(\pi(1+y^2))\,dy = \infty$, so $\phi$ is not in that domain (the brief
  \emph{Unphysicality of Hilbert spaces}, Proposition~3). This gives~(1). For~(2),
  $\psi \in D(U^{-1}YU)$ iff $U\psi \in D(Y)$, which fails.
\end{proof}

So the two observers do not describe the same $V$ with two families of defining observables, as
the premise requires: the second observer cannot express the first observer's $\psi$ in
$\mathcal S(\mathbb R)$. QSTO:V gives [B]WDQ only for changes of frame whose unitaries map $V$
onto itself. The restriction is stated in print by Carcassi, Calderón and Aidala: on the Schwartz
space ``the only unitary transformations we are interested in are those that map Schwartz
functions to Schwartz functions'', and the change of coordinate ``is automatically ruled
out''~\cite[Appendix, Sec.~VIII]{cca}. The brief does not repeat it. Theorem~\ref{thm:schwartz}
is the concrete instance, checked in Lean.

Which changes of variable do map $\mathcal S(\mathbb R)$ into itself is known: Galbis and Jordá
characterize them~\cite[Thm.~2.3]{galbis}. That the restriction cannot be lifted by choosing a
different space is shown in the companion note on the Hilbert brief~\cite[Sec.~2]{further-hilbert}:
no nonzero state of $L^2(\mathbb R)$ has finite position variance in every smooth frame both now
and after a free evolution. A quantum vector space has to come with its group of changes of frame
(Sec.~\ref{sec:framegroup}).

\subsection{Lean}

Module:\\ \lean{QuantumVectorSpaces.FurtherResults.SchwartzFrame},\\ namespace
\lean{QuantumVectorSpaces.FurtherResults}. The unitary $U$ (\lean{frameUnitary}), the packets
and the position operator are those of the Hilbert brief's formalization.
\begin{itemize}
  \item \lean{toLp_gaussSchwartz}: $\psi$ is $\pi^{-1/4}$ times TauCeti's Gaussian Schwartz map.
  \item \lean{toLp_mem_positionDomain}: every Schwartz function lies in the domain of the position
    operator on $L^2(\mathbb R)$.
  \item \lean{frameUnitary_gaussState_not_mem_range}, \lean{not_forall_frameUnitary_mem_range}:
    Theorem~\ref{thm:schwartz}(1).
  \item \lean{toLp_gaussSchwartz_not_mem_domain}: Theorem~\ref{thm:schwartz}(2).
\end{itemize}
The statement is made in $L^2(\mathbb R)$, where $U$ lives; it does not need $U$ to be smooth.

\section{The completeness problem of Definition 1}
\label{sec:irreducible}

The brief leaves one problem with its Definition~1 open (Sec.~5, ``Completeness of the Lie
algebra''). The definition admits two cases that the brief calls pathological: ``a Hilbert space
with the defining observables of a qubit'', which has ``too many vectors'', and ``a Schwartz space
with only position in the defining observable'', which has ``too few observables''. This section
shows that the second case is not a quantum vector space at all, and that the first one is
excluded by an irreducibility condition which Examples~3 and~4 of the brief satisfy.

Throughout, $V$ carries an inner product and defining observables $X_i$, and $\mathcal A$ is the
unital algebra they generate. The topology of the brief's Corollary~2 is given by the seminorms
$p_A(v) = \|Av\|$, $A \in \mathcal A$; Definition~1 requires $V$ to be complete and second
countable in it.

\subsection{Position alone}

\begin{theorem}\label{thm:posonly}
  Let $V = \mathcal S(\mathbb R)$ with the $L^2$ inner product, and let position $X$ be the only
  defining observable. Then $V$ is not complete in the topology of Corollary~2. So it is not a
  quantum vector space, whatever $\hbar$ and structure constants are chosen.
\end{theorem}

\begin{proof}
  Let $\psi_0(x) = \pi^{-1/4}e^{-x^2/2}$ and $f_m = u_m\psi_0$ with
  $u_m(x) = \sqrt{(m+1)^{-2} + x^2}$, a Schwartz function since $u_m$ is smooth with polynomially
  bounded derivatives. As $u_m(x) = |(m+1)^{-1} + ix|$, the triangle inequality gives
  $|u_m - u_k| \le |(m+1)^{-1} - (k+1)^{-1}|$ and $|u_m(x) - |x|| \le (m+1)^{-1}$. Every
  $A \in \mathcal A$ commutes with multiplication by $u_m$, so
  \[
    p_A(f_m - f_k) = \|(u_m - u_k)\,A\psi_0\| \le \bigl|(m+1)^{-1} - (k+1)^{-1}\bigr|\,\|A\psi_0\|,
  \]
  and $(f_m)$ is a Cauchy sequence. If it converged to some $h \in \mathcal S(\mathbb R)$, then
  $\|f_m - h\| \to 0$ (take $A = 1$), while $\|f_m - |x|\psi_0\| \le \|\psi_0\|/(m+1)$. So
  $h = |x|\psi_0$ almost everywhere, and everywhere since both are continuous. But $|x|\psi_0$ is
  not differentiable at~$0$.
\end{proof}

So the brief's second example is excluded by Definition~1 already, through completeness. The
pathology it was meant to show does not go away, though. The completion of $V$ is the space
$\Phi_X$ of all $f \in L^2(\mathbb R)$ with $x^k f \in L^2(\mathbb R)$ for every $k$, and $\Phi_X$
with position alone is a quantum vector space: it is a separable Fr\'echet space, the projective
limit of the weighted spaces $L^2((1+x^2)^k\,dx)$. Neither statement is formalized. It still has ``too few observables'': multiplication by
$e^{ix}$ commutes with position, on $\Phi_X$ as on $\mathcal S(\mathbb R)$.

\subsection{Irreducibility}

\begin{definition}\label{def:irreducible}
  $V$ with the defining observables $X_i$ is \emph{irreducible} if every linear map
  $T\colon V \to V$ that is continuous in the topology of Corollary~2 and commutes with every
  $X_i$ is a multiple of the identity.
\end{definition}

This is a form of the condition Roberts imposes on labelled Hilbert spaces: ``We simply suppose
that $\mathcal O' = \mathcal O_l'$, where $\mathcal O'$ denotes the set of all bounded operators
commuting with $\mathcal O$''~\cite[Sec.~2, p.~1099]{roberts}, where $\mathcal O$ is the set of all
observables and $\mathcal O_l$ that of the labelled ones. Without superselection rules
$\mathcal O' = \mathbb C 1$, and the condition says that the labelled observables have only the
scalars in their commutant. Definition~\ref{def:irreducible} states this on $V$ itself, with
continuous instead of bounded operators, and so needs no Hilbert completion. For unbounded
operators ``commutant'' can be made precise in several inequivalent ways; the statements below are
for this one. A sufficient condition for continuity is easy to check: if $T$ commutes with every
$X_i$, it commutes with every $A \in \mathcal A$, and $\|Tv\| \le C\|v\|$ for all $v$ gives
$p_A(Tv) = \|TAv\| \le C\,p_A(v)$.

\begin{theorem}\label{thm:irreducible}
  Let $\hbar \ne 0$.
  \begin{enumerate}
    \item Example~4 of the brief, $\mathbb C^2$ with the spin observables
      $S_i = \frac\hbar2\sigma_i$, is irreducible.
    \item $\mathbb C^2 \otimes \mathbb C^2$ with the defining observables $S_i \otimes 1$ is a
      quantum vector space, with the same commutation relations, but it is not irreducible:
      $1 \otimes \sigma_3$ commutes with every $S_i \otimes 1$.
    \item $\mathcal S(\mathbb R)$ with position alone is not irreducible: multiplication by
      $e^{ix}$ commutes with position.
    \item Example~3 of the brief, $\mathcal S(\mathbb R)$ with $X$, $P = -i\hbar\,d/dx$ and $I$,
      is irreducible.
  \end{enumerate}
\end{theorem}

\begin{proof}
  (1) A $2\times 2$ matrix that commutes with $\sigma_3$ is diagonal, and a diagonal one that
  commutes with $\sigma_1$ is scalar. (2) The commutation relations and symmetry pass from $S_i$
  to $S_i \otimes 1$, and the remaining conditions hold in every finite dimension; finite
  dimension also makes $1 \otimes \sigma_3$ continuous. (3) Multiplication by $e^{ix}$ is an
  isometry for the $L^2$ norm, hence continuous by the remark above, and it is not scalar.

  (4) The ladder operators $a_\mp = X \pm d/dx$ lie in $\mathcal A$, since
  $d/dx = (i/\hbar)P$. On the Hermite functions $\psi_n$ they act by
  $a_-\psi_n = \sqrt{2n}\,\psi_{n-1}$ and $a_+\psi_n = \sqrt{2(n+1)}\,\psi_{n+1}$. Write
  $c_n(f) = \langle\psi_n, f\rangle$ and $S_m f = \sum_{n<m} c_n(f)\psi_n$. The partial sums
  converge to $f$ in $\mathcal S(\mathbb R)$, because the coefficients of a Schwartz function
  decrease rapidly~\cite{simon}. Since $X$ is symmetric and $d/dx$ antisymmetric,
  $c_n(a_- g) = \langle a_+\psi_n, g\rangle = \sqrt{2(n+1)}\,c_{n+1}(g)$; so $a_- g = 0$ forces
  $c_n(g) = 0$ for $n \ge 1$, and $g = c_0(g)\psi_0$. Now let $T$ be continuous and commute with
  $X$ and $P$. It commutes with $a_\mp$, so $a_- T\psi_0 = T a_-\psi_0 = 0$ and
  $T\psi_0 = z\psi_0$ for some $z$. Applying $a_+$ gives $T\psi_n = z\psi_n$ for every $n$ by
  induction, hence $T S_m f = z S_m f$, and in the limit $Tf = zf$.
\end{proof}

Parts (1) and (2) also bear on the brief's claim, in the proof of its Theorem~5, that ``two
systems are physically equivalent if and only if their Lie algebras of observables are
equivalent''. The two spaces have isomorphic Lie algebras of defining observables and the same
structure constants, but they are not isomorphic, as their dimensions differ. Irreducibility
alone does not restore the claim: spin $\frac12$ and spin~$1$ are both irreducible, with the
same commutation relations. The system is determined by the representation, not by the Lie
algebra; for the Heisenberg algebra the Stone--von~Neumann theorem makes the irreducible
representation unique once $\hbar$ is fixed.

\subsection{The group form}

A stronger version of the same repair is to make a group part of the data. Let $G$ be a Lie
group, $\pi$ an irreducible unitary representation on a Hilbert space $H$, and let $V$ be the
space $H^\infty$ of smooth vectors, with the seminorms $\|d\pi(u)v\|$ for $u$ in the enveloping
algebra. These are the seminorms of Corollary~2 for the defining observables $i\,d\pi(Y_k)$, with
$(Y_k)$ a basis of the Lie algebra. Antoine builds his space from such a representation and
requires it to be invariant under the symmetry group~\cite{antoine2}; he considers $H^\infty$
and sets it aside because it need not be nuclear. For the Schr\"odinger representation of the
Heisenberg group, $H^\infty$ is $\mathcal S(\mathbb R)$ with its usual
topology~\cite[Secs.~1 and~6]{goodman}. In this form every defining observable is essentially
self-adjoint on $V$, as Antoine recalls from Segal and from Nelson and Stinespring, and it
generates a unitary group that maps $V$ onto itself. Every $\pi(g)$ maps $V$ onto itself, which
is the condition on changes of frame that Sec.~\ref{sec:schwartz} found missing from the brief's
argument for [B]WDQ. The cost
is that the definition presupposes the Hilbert space and the group. The Lean files contain only
the Lie-algebra form above; Theorem~\ref{thm:irreducible}(4) is its instance for the Heisenberg
group.

\subsection{Lean}

Module:\\ \lean{QuantumVectorSpaces.FurtherResults.Irreducibility},\\ namespace
\lean{QuantumVectorSpaces.FurtherResults}.
\begin{itemize}
  \item \lean{not_completeSpace_posOnly}, \lean{not_isQuantumVectorSpace_posOnly}:
    Theorem~\ref{thm:posonly}.
  \item \lean{IsIrreducible}: Definition~\ref{def:irreducible}; \lean{continuous_of_commute}:
    the sufficient condition for continuity.
  \item \lean{isIrreducible_qubit}, \lean{isQuantumVectorSpace_pair},
    \lean{not_isIrreducible_pair}, \lean{not_isIrreducible_posOnly},
    \lean{isIrreducible_schwartz}: Theorem~\ref{thm:irreducible}(1)--(4).
  \item The convergence of Hermite expansions in $\mathcal S(\mathbb R)$ is
    \lean{SchwartzMap.tendsto_hermiteSum}, in\\
    \lean{ForMathlib.Analysis.Distribution.SchwartzSpace.HermitePartialSum}.
\end{itemize}
Not formalized: that $\Phi_X$ is a quantum vector space, the remark on Theorem~5, and the group
form.

\section{[B]WDQ relative to a group of frames}
\label{sec:framegroup}

Section~\ref{sec:schwartz} and the companion note~\cite[Sec.~2]{further-hilbert} show that
[B]WDQ cannot hold for ``all frames'' as the brief words it: some changes of frame do not map
$V$ into itself, and no nonzero state serves every smooth frame at once. This section states
[B]WDQ with the frames as part of the data and shows that the brief's argument then goes through,
with a condition that can be checked.

\begin{definition}\label{def:framegroup}
  A \emph{frame group} for $V$ is a group $F$ with a homomorphism $\rho$ from $F$ to the linear
  bijections of $V$ such that every $\rho(g)$ preserves the inner product. The frame $g$ has the
  defining observables $X_i^g = \rho(g)^{-1}X_i\,\rho(g)$: the observer in frame $g$ describes the
  state $v$ by $\rho(g)v$ and measures $X_i$ on it. The frame group satisfies \emph{[B]WDQ} if for
  all $g, h \in F$ the defining observables of $g$ are observables of $h$, that is, symmetric
  operators continuous in the topology that Corollary~2 gives to the $X_i^h$.
\end{definition}

The last condition is the premise of the brief's proof of Theorem~5, for every pair of frames
in~$F$. The frame $1$ is the original one.

\begin{theorem}\label{thm:framegroup}
  Let the $X_i$ be symmetric and let $(F, \rho)$ be a frame group.
  \begin{enumerate}
    \item $(F, \rho)$ satisfies [B]WDQ if and only if every frame gives the same topology on $V$,
      and if and only if every $\rho(g)$ is continuous, that is, $F$ acts by automorphisms of the
      topological vector space~$V$.
    \item Then all frames have the same observables, and every $\rho(g)$ is a unitary equivalence
      of $V$ with itself. So it carries the pure states, with their topology, the Born
      probabilities and the observables of one frame to those of another.
    \item A sufficient condition is that every $X_i^g$ lies in the algebra $\mathcal A$ generated
      by the~$X_i$.
    \item The translations $f \mapsto f(\cdot - a)$, $a \in \mathbb R$, are a frame group of
      Example~3 that satisfies [B]WDQ, for every $\hbar$.
    \item The change of frame $y = \tan(\frac\pi2\operatorname{erf}x)$ belongs to no frame group of
      Example~3: no linear map of $\mathcal S(\mathbb R)$ agrees with its unitary $U$.
  \end{enumerate}
\end{theorem}

\begin{proof}
  (1) Since $\rho(g)$ preserves the inner product, the topology of the frame $g$ is the topology
  of the frame $1$ pulled back along $\rho(g)$; this is the transport result already used for
  isomorphisms (checked with the brief, \lean{QuantumVectorSpaces.induced_topology_eq_pullback}).
  Conjugation by $\rho(g)$ also preserves symmetry. By the brief's frame criterion
  (\lean{QuantumVectorSpaces.topology_eq_iff_continuous}), [B]WDQ for the pairs $(g, 1)$ and
  $(1, g)$ says that the topology of the frame $g$ equals that of the frame~$1$, and conversely
  equal topologies give [B]WDQ for every pair. A topology $t$ equals its pullback along
  $\rho(g)$ for every $g$ exactly when every $\rho(g)$ is continuous for $t$: continuity of
  $\rho(g)$ gives one inclusion, and continuity of $\rho(g^{-1}) = \rho(g)^{-1}$ gives the other.
  (2) follows from (1). (3) If every $X_i^g$ lies in $\mathcal A$, then
  $X_i = \rho(g)^{-1}X_i^{g^{-1}}\rho(g)$ lies in the algebra of the frame $g$, and two frames whose
  defining observables lie in each other's algebras give the same topology. (4) Under translation by
  $a$, position becomes $X + a$ and momentum stays $P$, both in $\mathcal A$; the inner product is
  preserved because Lebesgue measure is translation invariant. (5) is
  Theorem~\ref{thm:schwartz}(1).
\end{proof}

So the brief's argument for [B]WDQ is correct once ``all frames'' is read as ``all frames of a
given frame group that acts by automorphisms of $V$''. Part~(1) says that this restriction is
also necessary. Antoine requires the same invariance of his space under the symmetry
group~\cite[Sec.~II]{antoine2}, and his Proposition~4, on observables that transform by a
finite-dimensional representation, is related to part~(3). The negative half is not in his paper:
part~(5), and the companion note's result that no space works for all smooth frames. Which frame
group a quantum vector space should carry is a choice made with the space. For
$\mathcal S(\mathbb R)$ the changes of variable that preserve it are those of Galbis and
Jordá~\cite[Thm.~2.3]{galbis}, with the half-density factor added for unitarity; in the group form
of Sec.~\ref{sec:irreducible}, $F$ is the group $G$ itself.

\subsection{Lean}

Module:\\ \lean{QuantumVectorSpaces.FurtherResults.FrameGroups},\\ namespace
\lean{QuantumVectorSpaces.FurtherResults}.
\begin{itemize}
  \item \lean{FrameGroup}, \lean{frameObs}, \lean{IsFrameCompatible}:
    Definition~\ref{def:framegroup}.
  \item \lean{isFrameCompatible_iff_topology_eq}, \lean{isFrameCompatible_iff_continuous}:
    Theorem~\ref{thm:framegroup}(1).
  \item \lean{IsFrameCompatible.isObservable_iff}, \lean{IsFrameCompatible.unitaryEquiv}:
    Theorem~\ref{thm:framegroup}(2); the pure states and observables are carried by the brief's
    \lean{UnitaryEquiv.pureStates} and \lean{UnitaryEquiv.isObservable_conj_iff}.
  \item \lean{isFrameCompatible_of_mem_algebra}: Theorem~\ref{thm:framegroup}(3).
  \item \lean{translate}, \lean{frameObs_translate}, \lean{isFrameCompatible_translate}:
    Theorem~\ref{thm:framegroup}(4).
  \item \lean{not_exists_frameUnitary}: Theorem~\ref{thm:framegroup}(5).
\end{itemize}

\section{Ensembles over Example 3 and the continuity of entropy}
\label{sec:ensembles}

QSTO:V speaks of density operators and of the von Neumann entropy
$S(\rho) = -\operatorname{tr}(\rho\log\rho)$, and Sec.~5 of the brief calls the set of ensembles
``the most natural space to consider'', but the brief defines no ensemble space over~$V$. For
Example~3 this section proposes one and shows that the entropy is finite and continuous on it.
This is the quantum counterpart of the brief \emph{Entropy is not continuous in standard
probability spaces}~\cite{entropy}: there the entropy fails to be continuous in the $L^1$
topology, and here a topology made from the observables of $V$ restores continuity.

Take $\hbar = 1$, so that $N = \frac12(X^2 + P^2) - \frac12$ has the Hermite functions $\psi_n$
as eigenvectors, $N\psi_n = n\psi_n$. (For other $\hbar$ the Hermite functions are rescaled.) For a
positive trace-class operator $\rho$ on $L^2(\mathbb R)$ write
$\operatorname{tr}\rho N^j = \sum_n n^j\langle\psi_n, \rho\,\psi_n\rangle \in [0, \infty]$.

\begin{definition}\label{def:ensembles}
  A \emph{Schwartz density operator} is a positive operator $\rho$ on $L^2(\mathbb R)$ with
  $\operatorname{tr}\rho = 1$ and $\operatorname{tr}\rho N^j < \infty$ for every~$j$. The
  ensembles over Example~3 are the Schwartz density operators, with the topology of Schwartz
  operators of Keyl, Kiukas and Werner~\cite{kkw}.
\end{definition}

The definition agrees with theirs. They show that the seminorms
$\|H^m T H^{m'}\|_2$, $m, m' \in \mathbb N$, with $H = N + \frac12$, induce the topology of
Schwartz operators~\cite[Sec.~3.3]{kkw}. For positive $T$, the operators $H^m T H^{m'}$ are
products of the Hilbert--Schmidt operators $H^m T^{1/2}$ and $T^{1/2} H^{m'}$ exactly when every
$\operatorname{tr}(T H^{2m}) = \|T^{1/2}H^m\|_2^2$ is finite, that is when all moments
$\operatorname{tr}TN^j$ are finite. The eigenvectors of a Schwartz operator with nonzero eigenvalue are Schwartz
functions~\cite[Lemma~3.5(c)]{kkw}, so a Schwartz density operator is a mixture of pure states
of~$V$. The expectation of every polynomial in $X$ and $P$ is continuous on it, since
$|\operatorname{tr}(A\rho) - \operatorname{tr}(A\sigma)| \le \|A(\rho - \sigma)\|_1$ and the
right-hand side is a seminorm of the topology~\cite[eq.~(20)]{kkw}. In particular convergence
of Schwartz operators implies convergence in the trace norm and of $\operatorname{tr}\rho N$.

\begin{theorem}\label{thm:ensembles}
  \begin{enumerate}
    \item For every Schwartz function $f$ with $\|f\| = 1$, the pure state $|f\rangle\langle f|$
      is a Schwartz density operator.
    \item The entropy of a Schwartz density operator is finite.
    \item \emph{(Wehrl.)} Let $H = \sum_k \omega_k |b_k\rangle\langle b_k| \ge 0$ on a Hilbert
      space, with $\operatorname{tr}e^{-\beta H} < \infty$ for all $\beta > 0$. If density
      operators $\rho_m$ converge to $\sigma$ in the trace norm and
      $\operatorname{tr}\rho_m H \le h$ for all $m$, then $S(\rho_m) \to S(\sigma) < \infty$.
    \item Hence the entropy is continuous on the ensembles over Example~3: if
      $\|\rho_m - \sigma\|_1 \to 0$ and $\operatorname{tr}\rho_m N \to \operatorname{tr}\sigma N
      < \infty$, then $S(\rho_m) \to S(\sigma)$.
  \end{enumerate}
\end{theorem}

Part~(3) is Wehrl's~\cite[Sec.~II.D, p.~241]{wehrl}; that the condition on $H$ is also necessary
is Shirokov's~\cite[Prop.~1a]{shirokov}. Wehrl derives it from the lower semicontinuity of the
entropy and of the free energy. The proof below needs neither, nor the ordering of eigenvalues,
so that it could be checked in Lean with the trace-class theory at hand.

\begin{proof}
  (1) The Hermite coefficients $c_n = \langle\psi_n, f\rangle$ of a Schwartz function decrease
  rapidly~\cite{simon}: $(n+1)^{j+1}|c_n| \le C$. So
  $\operatorname{tr}(|f\rangle\langle f|N^j) = \sum_n n^j|c_n|^2 \le C^2\sum_n (n+1)^{-2}$.

  (3) Write $\rho = \sum_i \lambda_i|e_i\rangle\langle e_i|$ and $\eta(\lambda) = -\lambda\log\lambda$,
  so that $S(\rho) = \sum_i \eta(\lambda_i)$. For $K \in \mathbb N$ let
  $g_K(\lambda) = \min(-\log\lambda, K)$, continuous on $[0, 1]$, and split
  $\eta(\lambda) = \lambda g_K(\lambda) + r_K(\lambda)$ with $r_K \ge 0$.

  \emph{The head is continuous.} For every $n \ge 1$,
  $\rho^n - \sigma^n = \rho(\rho^{n-1} - \sigma^{n-1}) + (\rho - \sigma)\sigma^{n-1}$, so by
  induction $\|\rho^n - \sigma^n\|_1 \le n\|\rho - \sigma\|_1$ for density operators, and
  $|\operatorname{tr}\rho^n - \operatorname{tr}\sigma^n| \le n\|\rho - \sigma\|_1$. Given
  $\varepsilon > 0$, choose by Weierstrass a polynomial $q$ with $|g_K - q| \le \varepsilon$ on
  $[0, 1]$. Then $\sum_i\lambda_i q(\lambda_i) = \sum_j q_j\operatorname{tr}\rho^{j+1}$ is
  continuous in $\rho$, and it is within $\varepsilon\sum_i\lambda_i = \varepsilon$ of
  $\sum_i\lambda_i g_K(\lambda_i)$. So $\rho \mapsto \sum_i \lambda_i g_K(\lambda_i)$ is a
  uniform limit of continuous functions.

  \emph{The tail is uniformly small.} Klein's inequality
  $-a\log a \le -a\log b + b - a$ with $b = e^{-K-1-\beta\omega_k}$ gives
  $r_K(\lambda) \le \beta\omega_k\lambda + e^{-K-1}e^{-\beta\omega_k}$ for every~$k$ (for
  $\lambda \ge e^{-K}$ the left side is zero). Average over $k$ with the weights
  $|\langle b_k, e_i\rangle|^2$, which sum to one in $k$ and in $i$:
  \[
    \sum_i r_K(\lambda_i) \le \beta\operatorname{tr}\rho H + e^{-K-1}\operatorname{tr}e^{-\beta H}.
  \]
  On $\{\operatorname{tr}\rho H \le h\}$, choose $\beta$ with $\beta h < \varepsilon/4$, then $K$
  with $e^{-K-1}\operatorname{tr}e^{-\beta H} < \varepsilon/4$. The tails of $\rho_m$ and of
  $\sigma$ lie in $[0, \varepsilon/2)$, since the energy is lower semicontinuous and so
  $\operatorname{tr}\sigma H \le h$. Together with the head, $|S(\rho_m) - S(\sigma)| < \varepsilon$
  eventually. The same bound with $K = 0$ gives $S(\rho) < \infty$ whenever
  $\operatorname{tr}\rho H < \infty$, and so~(2), since $\operatorname{tr}e^{-\beta N} =
  (1 - e^{-\beta})^{-1}$.

  (4) If $\operatorname{tr}\rho_m N \to \operatorname{tr}\sigma N$, the energies are eventually
  bounded, and~(3) applies with $H = N$.
\end{proof}

The topology of Definition~\ref{def:ensembles} is stronger than needed: by~(4) the entropy is
already continuous for the trace norm together with $\operatorname{tr}\rho N$, and on each set
$\{\operatorname{tr}\rho N \le h\}$ for the trace norm alone. Whether the ensembles induce on the
pure states the quotient topology of~$V$, which the brief needs for its remark on second
countability, is not settled here. Woronowicz's description of the positive functionals on the
operators of $\mathcal S(\mathbb R)$ as the positive $\rho$ with $\rho L^2 \subset \mathcal S$
\cite{woronowicz71} suggests that the class of Definition~\ref{def:ensembles} is the natural one;
the comparison was not carried out.

\subsection{Lean}

The general theorem is in\\ \lean{ForPhyslib.QuantumMechanics.EntropyContinuity}, namespace
\lean{QuantumMechanics}; the application in\\
\lean{QuantumVectorSpaces.FurtherResults.Ensembles}, namespace
\lean{QuantumVectorSpaces.FurtherResults}.
\begin{itemize}
  \item \lean{numberMoment}, \lean{IsSchwartzDensity}: Definition~\ref{def:ensembles}, without the
    topology; the Hermite basis is TauCeti's \lean{hermiteHilbertBasis}.
  \item \lean{isSchwartzDensity_proj}: Theorem~\ref{thm:ensembles}(1).
  \item \lean{IsSchwartzDensity.vonNeumannEntropy_ne_top}, from \lean{vonNeumannEntropy_ne_top}:
    Theorem~\ref{thm:ensembles}(2).
  \item \lean{tendsto_vonNeumannEntropy}: Theorem~\ref{thm:ensembles}(3), with the energy bound
    required only eventually; \lean{traceENorm_pow_sub_pow_le}, \lean{tendsto_eigenSum} (the
    head), \lean{tailSum_le} (the tail), \lean{energy_le_of_tendsto} (lower semicontinuity of
    the energy).
  \item \lean{tendsto_vonNeumannEntropy_of_tendsto_numberMoment}: Theorem~\ref{thm:ensembles}(4).
\end{itemize}
Not formalized: the topology of Schwartz operators and the comparison with Keyl, Kiukas and
Werner's definition; the continuity of the expectations of polynomials in $X$ and $P$.

\begin{thebibliography}{9}
\bibitem{brief} G.~Carcassi, \emph{Physical vector spaces for quantum mechanics}, technical
  brief, Assumptions of Physics, \url{https://assumptionsofphysics.org/resources/briefs}.
\bibitem{hilbert} G.~Carcassi, \emph{Unphysicality of Hilbert spaces for quantum mechanics},
  technical brief, Assumptions of Physics,
  \url{https://assumptionsofphysics.org/resources/briefs}.
\bibitem{cca} G.~Carcassi, F.~Calder\'on, C.~A.~Aidala, The unphysicality of Hilbert spaces,
  Quantum Stud.: Math. Found. \textbf{12}, 13 (2025), arXiv:2308.06669v3.
\bibitem{galbis} A.~Galbis, E.~Jord\'a, Composition operators on the Schwartz space, Rev. Mat.
  Iberoam. \textbf{34} (2018) 397--412.
\bibitem{roberts} J.~E.~Roberts, The Dirac bra and ket formalism, J. Math. Phys. \textbf{7}
  (1966) 1097--1104.
\bibitem{antoine2} J.-P.~Antoine, Dirac formalism and symmetry problems in quantum mechanics.
  II. Symmetry problems, J. Math. Phys. \textbf{10} (1969) 2276--2290.
\bibitem{goodman} R.~Goodman, Analytic and entire vectors for representations of Lie groups,
  Trans. Amer. Math. Soc. \textbf{143} (1969) 55--76.
\bibitem{simon} B.~Simon, Distributions and their Hermite expansions, J. Math. Phys.
  \textbf{12} (1971) 140--148.
\bibitem{entropy} G.~Carcassi, \emph{Entropy is not continuous in standard probability spaces},
  technical brief, Assumptions of Physics,
  \url{https://assumptionsofphysics.org/resources/briefs}.
\bibitem{kkw} M.~Keyl, J.~Kiukas, R.~F.~Werner, Schwartz operators, Rev. Math. Phys.
  \textbf{28} (2016) 1630001, arXiv:1503.04086.
\bibitem{wehrl} A.~Wehrl, General properties of entropy, Rev. Mod. Phys. \textbf{50} (1978)
  221--260.
\bibitem{shirokov} M.~E.~Shirokov, Entropic characteristics of subsets of states I,
  arXiv:quant-ph/0510073 (2005).
\bibitem{woronowicz71} S.~L.~Woronowicz, The quantum problem of moments II, Rep. Math. Phys.
  \textbf{1} (1971) 175.
\bibitem{further-hilbert} R.~Stephan, \emph{Further results on the brief ``Unphysicality of
  Hilbert spaces for quantum mechanics''}, \lean{briefs/HilbertUnphysicality/further-results.pdf}.
\end{thebibliography}

\end{document}
